\section{Compatible sections and their mirror mixes}
Let $P$ be a chiral abstract $n$-polytope with a chosen base flag. We use the common paper's string rotation notation
\[
 \Gamma=U_n/N=\langle\sigma_1,\ldots,\sigma_{n-1}\rangle,
 \qquad \Lambda=\Gamma\diamond\overline\Gamma
                    \cong U_n/(N\cap\overline N).
\]
The symbol $\diamond$ means the subgroup of the direct product generated by the corresponding pairs of marked generators. It does not by itself assert polytopality. The common invariant-kernel theorem proves that $\Lambda$ is the rotation group of the smallest regular maniplex covering $P$. We denote this object by $R$. Its polytopality is the question to be tested.

Let $K,L,M$ be the facet, vertex-figure and medial section through the base flag. All their rotation groups are embedded in the same ambient $\Gamma$. Put
\[
 A=\Lambda_{n-1},\qquad B=\Lambda_0,\qquad C=\Lambda_{0,n-1},
 \qquad \pi:\Lambda\longrightarrow\overline\Gamma.
\]
The common local-mirror lemma gives
\begin{equation}\label{eq:sectionkernels}
 A\cap\ker\pi=X(K)\times1,\quad
 B\cap\ker\pi=X(L)\times1,\quad
 C\cap\ker\pi=X(M)\times1.
\end{equation}
The conjugation of the second coordinate in the vertex and medial cases is important: the global mirror is conjugation by $r_0$, whereas the intrinsic mirror of those sections is conjugation by $r_1$. Their relation $r_0=s_1r_1$ supplies exactly that conjugation. An abstract isomorphism between two chirality groups would not establish \eqref{eq:sectionkernels}.

We retain the standard string $C^+$ correspondence, with its recursive end-intersection criterion, as stated fully in the common paper. We also retain the Schulte--Weiss section theorem: facets of facets and vertex-figures of vertex-figures of a chiral polytope are directly regular. All sections of a directly regular polytope are directly regular. The next lemma proves the needed cover statement instead of retaining Cunningham's Proposition 2.2 as an input.

\begin{lemma}[A regular end]\label{lem:regularend}
If a chiral or directly regular polytope has directly regular facets or directly regular vertex-figures, its smallest regular cover is a polytope.
\end{lemma}
\begin{proof}
Induct on rank. Polygons have cyclic rotation group and are directly regular; the rank-one case is also immediate. A directly regular polytope is already its own smallest regular cover. Suppose now the facets $K$ are directly regular. Their mirror mixes are graphs of automorphisms, so $A$ is a marked copy of the facet rotation group, is polytopal, and satisfies $A\cap\ker\pi=1$. The facets of a vertex-figure $L$ are vertex-figures of the facets $K$, and hence are directly regular. The induction hypothesis therefore makes the mirror mix $B$ polytopal as well.

The second-coordinate marked group is a copy of the original mirrored polytope, so its intersection property gives
\[
                     \pi(A)\cap\pi(B)=\pi(C).
\]
The kernel condition in the common intersection-transport theorem holds because $A\cap\ker\pi=1$. Thus $A\cap B=C$. Both ends are polytopal, so the recursive string $C^+$ criterion makes $\Lambda$ polytopal. If the vertex-figures are directly regular instead, use the dual argument. This proves the induction without assuming the result to be established.
\end{proof}

For arbitrary chiral $P$, the section theorem now ensures that both $A$ and $B$ are polytopal. A chiral facet has directly regular facets, so Lemma~\ref{lem:regularend} applies to its mirror mix; the dual statement handles a chiral vertex-figure, and a directly regular end needs no induction. Consequently $R$ is a polytope exactly when $A\cap B=C$. This is the local prerequisite behind source Proposition 2.3.

\begin{theorem}[Cunningham's criterion]\label{thm:covercriterion}
For compatibly marked sections $K,L,M$ of a chiral polytope $P$,
\[
 R\text{ is a polytope}
       \quad\Longleftrightarrow\quad X(K)\cap X(L)\leq X(M).
\]
\end{theorem}
\begin{proof}
The image equality $\pi(A)\cap\pi(B)=\pi(C)$ holds as above. The common intersection-transport theorem turns $A\cap B=C$ into $A\cap B\cap\ker\pi\leq C$. Formula~\eqref{eq:sectionkernels} identifies this condition exactly with the displayed inclusion. The preceding paragraph supplies the polytopality of both ends, so the equality is equivalent to polytopality of $R$.
\end{proof}

\section{All general consequences in the source}
\begin{corollary}\label{cor:consequences}
The following statements hold with all subgroups embedded using the chosen base flag.
\begin{enumerate}[label=(\alph*)]
\item If either end section is directly regular, $R$ is polytopal.
\item If $M$ is directly regular, $R$ is polytopal exactly when $X(K)\cap X(L)=1$. Equivalently, this intersection meets $\Gamma_{0,n-1}$ trivially.
\item In rank four, $R$ is polytopal exactly when $X(K)\cap X(L)$ contains no nonidentity power of $\sigma_2$.
\item If $K$ and $L$ are totally chiral, then $R$ is polytopal exactly when $M$ is totally chiral.
\item If a chiral polytope $K$ satisfies $X(K)\cap\Gamma_0(K)=1$, then every chiral polytope having $K$ as facet has a polytopal smallest regular cover.
\end{enumerate}
\end{corollary}
\begin{proof}
A directly regular section has trivial chirality group, proving (a). In (b), $X(M)=1$; also
\[
 X(K)\cap X(L)\leq\Gamma_{n-1}\cap\Gamma_0
                            =\Gamma_{0,n-1},
\]
so taking a further intersection with the medial group does not change the condition. For (c), the medial section is a polygon, its rotation group is $\langle\sigma_2\rangle$, and its chirality group is trivial. Every nontrivial element in the displayed intersection is therefore a nonidentity power of $\sigma_2$. The word ``nonidentity'' is necessary: the identity is in every subgroup.

For (d), total chirality means $X(K)=\Gamma_{n-1}$ and $X(L)=\Gamma_0$. Their intersection is the full medial group, so the criterion holds exactly when $X(M)=\Gamma_{0,n-1}$. Finally, in (e) the common medial subgroup is the subgroup $\Gamma_0(K)$ inside the facet. Hence $X(K)\cap X(L)$ is contained in the assumed trivial intersection, proving the criterion for every such extension. This proves no existence statement for extensions of a prescribed rank.
\end{proof}

These are respectively source Remark 3.2, Corollaries 3.3--3.4, Proposition 3.6 and Corollary 3.7. The criterion itself is source Theorem 3.1. The first-coordinate identity for chirality in the common paper recovers source Proposition 2.4; the argument above also recovers Propositions 2.2--2.3.
